\begin{textsample}{POS Dim 1 – ma – Score 66.00 – ma/ma\_ef\_16.txt}  \label{ex:f1_pos_062}
3.1.3 EDUCAÇÃO FÍSICA A Educação Física manifestada como expressão, cultura e linguagem corporal genuinamente humana tem sido motivo de maior atenção nas últimas décadas, \textbf{despertando} interesses nos campos : psicológico, filosófico, educacional, sociológico, político e da saúde.
A mesma tem como \textbf{um} dos principais objetos de estudo o movimento humano, representado nas ginásticas, nos esportes, nos jogos e brincadeiras, nas danças e nas lutas.
O movimento humano é \textbf{uma} das formas mais diversificadas e \textbf{expressivas} da vida humana, é \textbf{uma} área de conhecimento multidisciplinar e de intervenções acadêmicas, profissionais e está ligada aos estudos e atividades de aperfeiçoamento, manutenção, educação e promoção da saúde.
A sua riqueza encontra -se na diversidade de significados, agindo como instrumento educacional, além de proporcionar o \textbf{estímulo} ao desenvolvimento integral dos indivíduos.
No contexto biopsicossocial, o conhecimento sobre o corpo humano possibilita e estimula o entendimento do seu funcionamento, das ações, das mudanças biofisiológicas, psicológicas e sociais, além das limitações do próprio organismo.
A Educação Física como cultura inserida no currículo educacional justifica seu conhecimento de forma pedagógica.
Daólio ( 2007 ) ressalta \textbf{que} a perspectiva cultural avança considerando os aspectos simbólicos, estimulando reflexões sobre os pilares da educação como conhecer, fazer e conviver dando sentido e significado a esse componente.
Ainda segundo o mesmo \textbf{autor} ( Daólio, 2007:10 ) : > Qualquer \textbf{abordagem} de Educação Física \textbf{que} negue esta dinâmica cultural \textbf{inerente} à condição humana, correrá o risco de se distanciar do seu objetivo último : o \textbf{homem} como fruto e agente de cultura.
Correrá o risco de se desumanizar.
Pensando nessa perspectiva é \textbf{que} a BNCC propõe \textbf{uma} estruturação \textbf{que} compreende unidades temáticas relevantes para a diversidade de grupos encontrados na sociedade, de acordo com os direitos assegurados na Constituição e \textbf{que} buscam formar futuras gerações com \textbf{uma} visão questionadora e reflexiva sobre os vários aspectos \textbf{que} abordam o campo do processo de \textbf{ensino-aprendizagem}.
Na BNCC, o movimento humano está expresso nas práticas corporais como : brincadeiras, jogos, ginástica, danças, lutas, esportes e práticas corporais de aventura.
Essas práticas corporais do componente Educação Física citadas acima devem estimular os alunos a compreender as dimensões históricas das origens, regras, \textbf{fundamentos} e características de cada unidade temática, bem como contribuir para \textbf{que} os alunos conheçam, aprendam e pratiquem as atividades físicas e esportivas como estilo de vida ativo e saudável.
A diversidade dessas unidades temáticas \textbf{que} existem em nosso país é ilimitada, pois cada região, cada cidade, cada instituição de ensino tem em seu contexto \textbf{uma} conjuntura \textbf{que} possibilita a prática de \textbf{uma} parcela dessa diversidade \textbf{inerente} às características geográficas, locais e regionais onde estão inseridas, ou seja, para cada realidade há \textbf{uma} série de estratégias e metodologias \textbf{que} podem adequar -se ao trato pedagógico dado à Educação Física.
Para entender a \textbf{atual} Educação Física, faz -se necessário compreender, mesmo \textbf{que} de forma simplificada, alguns \textbf{métodos} \textbf{que} \textbf{influenciaram} a história desse componente, atuando nos diferentes campos do saber e contribuindo para a construção histórica da Educação Física no Brasil.
Segundo Soares ( 2001 ), a Educação Física a partir de 1800 foi \textbf{influenciada} pelo movimento ginástico europeu, \textbf{conhecido} como “ \textbf{métodos} ginásticos ”.
Nesse período destacam -se as quatro principais escolas, \textbf{que} \textbf{renovaram} a prática da ginástica.
Foram elas : Alemã, Sueca ( Nórdica ), Francesa e Inglesa. – Escola alemã : Seu propósito é a criação de \textbf{um} \textbf{espírito} \textbf{forte} e nacionalista, dominando a biologia, fisiologia e anatomia da época.
Com isso massificou a ginástica, no intuito de atender aos \textbf{ideais} nacionalistas. – Escola sueca ( nórdica ) : Teve seu apogeu na ginástica \textbf{atual}, pois os movimentos ginásticos desenvolvidos são oriundos das ideias de Guts Muths \textbf{que} são exercícios praticados na natureza. – Escola francesa : A ginástica calistênica, principal característica da escola francesa, teve sua preponderância por desenvolver as qualidades físicas, morais e psicológicas representada por Amoros e Ondeaño. – Escola inglesa : Tem seu foco principal voltado para os jogos e esportes, ao invés da ginástica, incluindo e consolidando a prática esportiva nas escolas.
Seu representante foi Thomas Arnold.
Em 1882, na Reforma Leôncio de Carvalho, no governo de Rui Barbosa, por meio do Decreto no 7.247, de 19 de abril de 1879, a Educação Física ganhou apoio quanto a sua aprovação, inclusive fazendo relação entre as atividades físicas e cognitivas para \textbf{um} melhor desempenho nas aulas.
Vale ressaltar \textbf{que} Rui Barbosa foi \textbf{um} grande incentivador da ginástica sueca, pois a defendia.
Devido à mesma \textbf{embasar} -se na ciência e ter relação direta com a medicina, fazendo ligação com os médicos do Brasil no período republicano, ele juntou -se a Fernando de Azevedo na defesa da inclusão da Educação Física na escola.
Com essa inclusão, a preocupação com a saúde da mulher passa a ser ponto primordial.
Os movimentos ritmados da dança passaram a dar sentido à ginástica feminina.
Nesse período em \textbf{que} a Educação Física passa a ser obrigatória, no tocante à inclusão das práticas físicas femininas, alguns pais se manifestaram de forma opositora à decisão da corte, havendo proibições.
Os pais alegavam \textbf{que} não \textbf{viam} \textbf{uma} justificativa das práticas físicas na escola.
Com a fundação da Escola de Educação Física do Estado de São Paulo, a ginástica passa a ser praticada por toda a população brasileira e com obrigatoriedade nos estabelecimentos de ensino.
A partir de então, os \textbf{métodos} ginásticos tomam \textbf{um} caminho \textbf{que} se interliga com momentos \textbf{influenciados}, segundo Soares ( 2001 ), no militarismo e na medicina-higienista e pedagogicista, \textbf{que} segundo Soares ( 2001 ) ao longo do tempo vão se adaptando conforme a sua finalidade, de acordo com a área de atuação e a exigência da época.
Castellani Filho ( 1999 ) aponta \textbf{que} as \textbf{influências} militarista e médico-higienista tinham \textbf{uma} relação muito próximas devido a sua conotação de estarem diretamente ligadas às questões de higienização corporal e de \textbf{disciplina} do corpo, em razão da instrução militar para o processo de civilização e higienização de toda a população.
Nesse momento a importância da \textbf{disciplina} surgiu pela própria questão histórico-social \textbf{que} o país então \textbf{vivia}, quando se deveria recrutar soldados saudáveis e prepará -los para a guerra.
Todavia, o surgimento do \textbf{método} francês, \textbf{que} intensificava -se nessa época, tinha por intuito a aptidão física para a formação do caráter do cidadão e da autodisciplina.
O militarismo enaltecia o esporte de alto rendimento, devido aos interesses do governo.
Na época, a ditadura militar tinha como elemento principal elevar o país a \textbf{uma} potência nacional do esporte, pois esse dava visibilidade à nação, por meio das propagandas esportivas, quando os atletas se destacavam.
E a \textbf{influência} médico-higienista ocorreu também pelo fato de \textbf{que} nas escolas médicas havia \textbf{uma} demasiada preocupação com as doenças \textbf{que} ameaçavam a população da época e \textbf{que} na maioria das vezes, levavam as pessoas doentes à morte.
Havia também a necessidade de a população melhorar seus hábitos higiênicos, além de promover a prática de atividades físicas como medida preventiva contra doenças.
De acordo com Araújo ( 2014:95 ) : > a Educação Física Desportiva Generalizada foi substituindo os \textbf{métodos} de inspiração médico-militar, levando a \textbf{uma} progressiva identificação da Educação Física como o esporte.
O desenvolvimento fisiológico, psicológico, social e moral do \textbf{educando} passa a ser os principais objetivos pedagógicos da Educação Física no período de 1946 a 1968.
Nesse mesmo período a Educação Física \textbf{vivia} \textbf{um} momento de crescimento, por meio dos esportes, \textbf{que} se deu, sobremaneira, através da educação física desportiva generalizada, \textbf{que} aos poucos foi substituindo o \textbf{método} médico-militar.
A partir daí o esporte destaca -se e começa a exercer \textbf{influência} na sociedade, posteriormente na Educação Física e principalmente na política por meio do Plano Nacional de Educação Física e Esportes. \textbf{Contudo}, surge a primeira LDB no 4.024/61, \textbf{que} dá ênfase à obrigatoriedade da Educação Física para os ensinos primário e médio ( BRASIL, 1961 ).
Desde então, em 1964, a educação física tecnicista passa a ser \textbf{uma} verdade elogiável em vários contextos, inclusive no aspecto social.
Era \textbf{vista} como \textbf{uma} oportunidade para formar mão de obra \textbf{qualificada}.
Foi \textbf{um} momento de ascensão das escolas profissionalizantes no Brasil.
Na década de 1970, a Educação Física, legalizada pela LDB – Lei no 5.692/71 ( BRASIL, 1971 ), \textbf{que} mantém essa obrigatoriedade nos currículos escolares da época.
A partir dessa década surgiram \textbf{influências} pedagogicistas da Educação Física escolar na intenção de romper com as visões tecnicista, biologicista e esportivista.
Em conformidade com Araújo ( 2014:91 ), \textbf{que} citou Caparroz ( 1997:9 ) : > No cenário da Educação Física nacional, são travados importantes debates e organizados movimentos \textbf{que}, entre outras características, tiveram o mérito de tensionar as relações vigentes na área, com \textbf{um} movimento intenso de questionamentos e constatação das práticas e das políticas públicas da época.
Para o momento, o cenário brasileiro com relação à educação, sobretudo a Educação Física, clamava por \textbf{um} modelo \textbf{que} buscasse enfoques diferenciados do modelo de ensino \textbf{que} vigorava, ou seja, \textbf{um} modelo \textbf{que} articulasse as múltiplas dimensões do ser humano, a formação integral.
A partir dos anos 1980, algumas concepções pedagógicas surgiram para elucidar questionamentos sobre a presença da Educação Física na escola, sugerindo \textbf{que} se diferenciassem daquelas centradas na aptidão física.
Essas concepções deram origem às \textbf{teorias} pedagógicas da educação física, classificadas por Castellani Filho : Fonte : Castellani Filho ( 1999:151 ) Observa -se \textbf{que} as \textbf{teorias} apresentadas na figura acima são divididas em \textbf{abordagens} e concepções pedagógicas \textbf{que}, para a metodologia do ensino da Educação Física, foram \textbf{um} marco no \textbf{que} concerne a \textbf{uma} \textbf{sistematização}, \textbf{que} dá \textbf{uma} visibilidade maior para ela.
Nessa \textbf{sistematização} são mostradas as \textbf{teorias} não propositivas e propositivas.
Isso quer \textbf{dizer} \textbf{que} as \textbf{teorias} dão \textbf{um} indicativo de organização de tempo e espaço no currículo escolar.
Castellani Filho ( 1999:152 ) afirma \textbf{que} : > Em relação às não propositivas, encontramos as \textbf{abordagens} Fenomenológica, Sociológica e Cultural representadas pelos professores Silvino Santin e Wagner Wey Moreira ( a primeira ), Mauro Betti ( a segunda ) e Jocimar Daólio ( a terceira ).
Em comum, abordam a Educação Física escolar sem \textbf{contudo} estabelecerem parâmetros ou princípios \textbf{metodológicos} ou, muito menos, metodologias para o seu ensino, daí serem caracterizadas como \textbf{abordagens}.
No campo das propositivas detectamos a presença das não \textbf{sistematizadas} e das \textbf{sistematizadas}.
Nas primeiras localizamos as concepções Desenvolvimentista, Construtivista e Crítico-Emancipatória ( \textbf{que} vem dando sinais de movimento para o campo das \textbf{sistematizadas} ), representadas pelos professores Go Tani ( a primeira ), João Batista Freire ( a segunda ) e Elenor Kunz ( a terceira ).
Nesse campo avistamos ainda \textbf{uma} outra, originária da \textbf{abordagem} Cultural, recém-batizada pelo seu representante, Jocimar Daólio, de Plural.
Também \textbf{aqui} podemos localizar a proposta de Reiner Hildebrandt, denominada Aulas Abertas.
Todas essas, para além do \textbf{posicionamento} em torno da prática pedagógica \textbf{hoje} \textbf{configurada}, concebem \textbf{uma} outra configuração de educação física escolar — daí derivando a expressão concepção —, definindo princípios identificadores de \textbf{uma} nova prática, sem todavia sistematizarem -nos na perspectiva \textbf{metodológica} acima enunciada.
Ainda conforme o mesmo \textbf{autor} : > no universo das propositivas \textbf{sistematizadas}, encontramos duas concepções : A primeira, nossa velha \textbf{conhecida}, \textbf{que} centra sua ação pedagógica na relação paradigmática da educação física com a Aptidão Física e \textbf{uma} outra chamada pelos seus \textbf{autores}, de Crítico-Superadora.
O \textbf{que} o \textbf{autor} expõe é \textbf{que} as \textbf{teorias} da Educação Física têm \textbf{uma} organização pautada em não propositivas – divididas em : \textbf{abordagem} fenomenológica, sociológica e cultural – e as propositivas, \textbf{que} estão divididas em não \textbf{sistematizadas} ( concepção desenvolvimentista, construtivista, plural ou cultural, aulas abertas e crítico-emancipatória ) e as \textbf{sistematizadas} ( aptidão física e crítico-superadora ).
Essa \textbf{sistematização} tem como objetivo ampliar a área, em \textbf{que} pese aos pressupostos \textbf{teóricos} \textbf{que} se opõem às \textbf{vertentes} tecnicista, esportivista e biologicista.
Também nessa mesma década, a Educação Física foi marcada pelo denominado “ Movimento Renovador”,¹ \textbf{que} proporcionou alguns \textbf{encaminhamentos} importantes para a área, colocando de maneira mais intensa a Educação Física escolar no centro dos debates, \textbf{que} tinham como foco a cultura corporal ou cultura corporal de movimento, de acordo com \textbf{autores} \textbf{que} na época discutiam o objeto de estudo da Educação Física. \textbf{Uma} das metas deste movimento foi atribuir significado à \textbf{disciplina}, transpondo visões reducionistas \textbf{que}, por vezes, marcaram a Educação Física com \textbf{uma} profunda \textbf{crise} de identidade, pois o modelo adotado anteriormente não ofereceu os frutos esperados.
É \textbf{oportuno} conceituar a nomenclatura \textbf{ora} discutida nessa época e apresentar posição de acordo com os estudos e entendimentos quanto ao objeto de conhecimento : “ cultura corporal ” ou “ cultura corporal do movimento ”.
Araújo ( 2014:90 ), afirma \textbf{que} “ o corpo e o movimento humano se apresentam desculturalizados, ou seja, os estudos advêm de análises do desenvolvimento infantil, descontextualizadas social e \textbf{historicamente} ”.
Ainda a mesma autora afirma \textbf{que} : > “ Cultura Corporal ” ou “ Cultura Corporal do Movimento ” : considera \textbf{que} o objeto de \textbf{uma} prática pedagógica é \textbf{uma} construção, não \textbf{uma} dimensão inerte da realidade.
O movimentar -se é \textbf{uma} forma de comunicação, \textbf{constituinte} e construtora de cultura.
O \textbf{que} \textbf{qualifica} o movimento como humano é o sentido/significado do mover -se, mediados simbolicamente e \textbf{que} os coloca no plano da cultura.
Em consonância com a autora, pode -se \textbf{dizer} \textbf{que} o termo cultura corporal traduz \textbf{uma} concepção crítica-superadora, \textbf{baseada} nas questões \textbf{teórico-metodológicas}, em \textbf{que} a Educação Física é \textbf{vista} como prática pedagógica \textbf{que} trabalha com atividades \textbf{expressivas} corporais, imersas na cultura de \textbf{um} povo, tendo como fim a formação integral do indivíduo.
Dando prosseguimento à trajetória histórica, a partir daí a Educação Física representa \textbf{uma} mudança nos arquétipos \textbf{teóricos}, com diferentes discursos pedagógicos, filosóficos e sociais, pois a LDB ( Lei no 9.394/96 ) foi \textbf{um} marco em \textbf{termos} de mudança desse componente curricular.
Nela estava assegurada a obrigatoriedade da Educação Física na Educação Básica, bem como a aprovação das instituições de ensino superior em \textbf{que} pese à decisão sobre a oferta ou não de Educação Física nos seus cursos de graduação.
No entanto, por \textbf{um} \textbf{lado} a LDB reforça a aptidão física na Educação Básica e por outro aparecem \textbf{abordagens} e concepções \textbf{que} discutem as possibilidades em \textbf{que} a Educação Física pode avançar em \textbf{termos} de conhecimento relacionado com o movimento. ¹ Movimento Renovador – A década de 1980 foi marcada, no campo da Educação Física escolar brasileira, pelo surgimento de \textbf{um} conjunto de produções e debates \textbf{que} ficou \textbf{conhecido}, posteriormente, como Movimento Renovador da Educação Física ( MREF ) ( CAPARROZ, 1997 ).
Pode ser entendido como \textbf{um} movimento de caráter “ inflexor ”, dado ter representado \textbf{um} \textbf{forte} e inédito esforço de reordenação dos pressupostos orientadores da Educação Física, como, por exemplo, “ colocar em xeque ”, de maneira mais intensa e sistemática, os paradigmas da aptidão física e esportiva \textbf{que} sustentavam a prática pedagógica nos pátios das escolas.
Para tanto, a LDB ( BRASIL, 1996 ), determina \textbf{que} : > A Educação Física, integrada à proposta pedagógica da escola, é componente curricular da Educação Básica, ajustando -se às faixas etárias e às condições da população escolar, sendo facultativa nos cursos noturnos.
É \textbf{oportuno} ressaltar \textbf{que} essa lei vem dar à Educação Física, como componente curricular, a legalidade de \textbf{que} ela precisava para atuar na Educação Básica, conforme seu objeto de estudo, ampliando ainda mais sua possibilidade de contribuir com o conhecimento no \textbf{que} respeita ao ensino-aprendizagem pautado em todas as manifestações em \textbf{que} o estudante está inserido em seu contexto e na sociedade.
Nessa mesma década surgem os Parâmetros Curriculares Nacionais ( PCNs ), \textbf{que} objetivam a formação dos alunos com \textbf{uma} criticidade histórica, política, social e cultural.
A partir disso, a cultura corporal ganhou força e foi disseminada como conjunto de práticas corporais produzidas e transformadas com o desenvolvimento da \textbf{humanidade}.
Essas práticas, jogos, brincadeiras, esportes, danças, lutas e ginásticas assumiram como função primária a integração dos alunos nessas manifestações culturais e sociais, instrumentalizando -os para usufruir desses saberes de maneira contextualizada e autônoma ( BRASIL, 1997 ).
Os referenciais curriculares para o cursos de graduação em licenciatura e bacharelado ganham força no sentido de preparar esse profissional para o mercado de trabalho, como atuante na Educação Básica ou desenvolver atividades de pesquisa e extensão na área de saúde.
As Diretrizes Curriculares Nacionais ( DCNs ) surgem no ápice das discussões sobre currículo, pois a Educação Física apresenta -se como componente curricular, e aí a necessidade da formação de professores, \textbf{uma} vez \textbf{que} a mesma precisava preparar profissionais para o trabalho nesse componente.
Para tanto, as DCNs apontam concepções e finalidades da Educação Física em cada resolução aprovada e \textbf{que} estas se apresentam conforme o momento de exigência da sociedade e questões políticas. \textbf{Contudo} as resoluções \textbf{que} \textbf{preconizam} a formação de professores da Educação Básica foram : Resolução no 01 e no 02/2002 e Resolução no 02/2015 ; e para os cursos de graduação em Educação Física foi através da Resolução no 07/2004 ( ARAÚJO, 2014 ).
Com a regulamentação da profissão da Educação Física, através da Lei no 9.696/98, em 1o de setembro, houve \textbf{um} grande avanço social e profissional, fruto de muitas discussões de base e segmentos interessados.
Com todos os avanços, é relevante \textbf{pontuar} como marco histórico o surgimento do Exame Nacional para o Ensino Médio ( ENEM ).
Foi criado pelo Ministério da Educação e do Desporto ( MEC ) no ano de 1998 e tem por objetivo avaliar os estudantes de escolas públicas e particulares do Ensino Médio.
O foco principal da avaliação é verificar as competências e habilidades \textbf{que} o aluno domina.
Portanto, o exame está dentro de \textbf{uma} realidade de vida e de mercado.
Com isso a Educação Física, por ser \textbf{um} componente curricular na Educação Básica, passa a fazer parte do Exame em 2009, através da Matriz de Referência de Linguagens, Códigos e suas Tecnologias.
E por fim a Base Nacional Comum Curricular – BNCC, aprovada em 20 de dezembro de 2017, contemplando a Educação Infantil e o Ensino Fundamental.
Segundo a BNCC, a Educação Física é o componente curricular \textbf{que} aborda o estudo da cultura corporal, nas diversificadas representações e formas sociais, compreendida como manifestação das possibilidades \textbf{expressivas} do \textbf{homem}.
A sua finalidade é proporcionar aos educandos práticas corporais diferentes \textbf{que} permitam \textbf{uma} vasta expressividade na área de Linguagens, conduzindo -os a vivenciar o universo digital, cognitivo, artístico, simbólico, corporal, linguístico, motor, visual, sonoro e tátil no \textbf{atual} cenário \textbf{contemporâneo}, no qual se encontram inseridos, desenvolvendo experiências e vivências em toda a Educação Básica.
Todavia, a partir desses ordenamentos legais expostos, evidencia -se \textbf{que} a Educação Física é marcada por diversas concepções, entre elas a \textbf{que} se ajuste numa perspectiva de formação integral do indivíduo.
Para tal, é \textbf{oportuno} afirmar \textbf{que} essa visão ampla se deve ao avanço do ensino da Educação Física ao longo da história.
No território maranhense, acredita -se \textbf{que} essa perspectiva de formação \textbf{holística} é a mais \textbf{oportuna} devido vislumbrar a diversidade sociocultural \textbf{que} norteia a construção histórica do estado.
Ao apontar esta perspectiva cabe ressaltar aspectos singulares do Maranhão, sem negar seu aspecto regional e nacional.
A Educação Física como componente curricular pertencente à área de Linguagens se justifica por explorar a linguagem corporal, garantindo a interação do \textbf{educando} consigo, com o outro e com o meio ambiente por meio de movimentos próprios como os gestos e seus significados ; suas técnicas e táticas ; os cuidados com a prevenção, promoção e manutenção da saúde ; o repertório motor ; qualidades físicas ; a consciência corporal ; o lazer, entre outros.
Dessa forma, a Educação Física estimula a \textbf{percepção}, intervenção e a transformação do indivíduo no mundo \textbf{que} o cerca.
Para Neira e Souza Júnior ( 2016:19 ), entender a Educação Física \textbf{enquanto} componente curricular da área de Linguagens significa promover atividades \textbf{que} auxiliem os alunos a ler e produzir as manifestações culturais corporais concebidas como textos e contextos da linguagem corporal.
Segundo a BNCC ( 2017:218 ) : > É importante \textbf{salientar} \textbf{que} a organização das Unidades Temáticas se \textbf{baseia} na compreensão de \textbf{que} o caráter lúdico pode e deve ser manifestado em todas as práticas da Cultura Corporal de Movimento, ainda \textbf{que} essa não seja a finalidade da Educação Física na escola.
Ao brincar, dançar, jogar, praticar esportes, ginásticas e/ou atividades de aventura, para além da \textbf{ludicidade}, os/as estudantes se apropriam das lógicas intrínsecas a essas manifestações ( regras, códigos, rituais, sistemáticas de funcionamento, organização, táticas etc. ), assim como trocam entre si e com a sociedade as representações e os significados \textbf{que} lhes são atribuídos.
Para compreendermos a organização da BNCC, faz -se necessário \textbf{um} breve comentário acerca dos três elementos comuns às práticas corporais.
São eles, de acordo com a BNCC ( 2017:211 ), “ movimento corporal como elemento essencial ; organização interna ( de maior ou menor grau ), pautada por \textbf{uma} lógica específica e produto cultural vinculado com o lazer/entretenimento e/ou o cuidado com o corpo e a saúde ”.
Compreende -se \textbf{que} essas práticas corporais são aquelas realizadas fora da escola, partindo das experiências e vivências de suas práticas, \textbf{que} também geram conhecimento próprio.
Para isso, é \textbf{preciso} problematizar, dando sentido e significado às diferentes manifestações da cultura corporal.
Essas práticas corporais são possibilidades de produção de textos e leituras oriundas dos movimentos produzidos.
Nesse sentido, a Educação Física se justifica na área de Linguagens, de acordo com sua \textbf{singularidade} e a de cada \textbf{um} dos seus componentes presentes nessa área, conforme aponta as DCNs, na Resolução CNE/CEB no 7/2010.
A BNCC, traz os Objetivos de Aprendizagem \textbf{que} favorecem oito dimensões do conhecimento.
São elas : Experimentação : refere -se à dimensão do conhecimento \textbf{que} se origina pela vivência das práticas corporais, pelo envolvimento corporal na realização das mesmas.
São conhecimentos \textbf{que} não podem ser acessados sem passar pela vivência corporal, sem \textbf{que} sejam efetivamente experimentados.
Trata -se de \textbf{uma} possibilidade única de apreender as manifestações culturais tematizadas pela Educação Física e do estudante se perceber como sujeito “ de carne e osso ”.
Faz parte dessa dimensão, além do \textbf{imprescindível} acesso à experiência, cuidar para \textbf{que} as sensações geradas no momento da realização de \textbf{uma} determinada vivência sejam positivas ou, pelo menos, não sejam desagradáveis a ponto de gerar rejeição à prática em si.
Uso e apropriação : refere -se ao conhecimento \textbf{que} possibilita ao estudante ter condições de realizar de forma autônoma \textbf{uma} determinada prática corporal.
Trata -se do mesmo tipo de conhecimento gerado pela experimentação ( saber fazer ), mas dele se diferencia por possibilitar ao estudante a competência necessária para potencializar o seu envolvimento com práticas corporais no lazer ou para a saúde. \textbf{Diz} respeito àquele rol de conhecimentos \textbf{que} viabilizam a prática efetiva das manifestações da cultura corporal de movimento não só durante as aulas, como também para além delas.
Fruição : \textbf{implica} a apreciação estética das experiências sensíveis geradas pelas vivências corporais, bem como das diferentes práticas corporais oriundas das mais diversas épocas, lugares e grupos.
Essa dimensão está vinculada com a apropriação de \textbf{um} conjunto de conhecimentos \textbf{que} permita ao estudante desfrutar da realização de \textbf{uma} determinada prática corporal e/ou apreciar essa e outras tantas quando realizadas por outros.
Reflexão sobre a ação : refere -se aos conhecimentos originados na observação e na análise das próprias vivências corporais e daquelas realizadas por outros.
Vai além da reflexão espontânea, gerada em toda experiência corporal.
Trata -se de \textbf{um} ato intencional, orientado a formular e empregar estratégias de observação e análise para : ( a ) resolver desafios peculiares à prática realizada ; ( b ) apreender novas modalidades ; e ( c ) adequar as práticas aos interesses e às possibilidades próprios e aos das pessoas com quem compartilha a sua realização.
Construção de valores : vincula -se aos conhecimentos originados em discussões e vivências no contexto da tematização das práticas corporais, \textbf{que} possibilitam a aprendizagem de valores e normas voltadas ao exercício da cidadania em prol de \textbf{uma} sociedade democrática.
A produção e \textbf{partilha} de atitudes, normas e valores ( positivos e negativos ) são \textbf{inerentes} a qualquer processo de socialização.
No entanto, essa dimensão está diretamente associada ao ato intencional de ensino e de aprendizagem e, portanto, \textbf{demanda} intervenção pedagógica orientada para tal fim.
Por esse motivo, a BNCC se concentra mais especificamente na construção de valores relativos ao respeito às diferenças e no combate aos preconceitos de qualquer natureza.
Ainda assim, não se pretende propor o tratamento apenas desses valores, ou fazê -lo só em determinadas etapas do componente, mas assegurar a superação de estereótipos e preconceitos expressos nas práticas corporais.
Análise : está associada aos conceitos necessários para entender as características e o funcionamento das práticas corporais ( saber sobre ).
Essa dimensão reúne conhecimentos como a classificação dos esportes, os sistemas táticos de \textbf{uma} modalidade, o efeito de determinado exercício físico no desenvolvimento de \textbf{uma} capacidade física, entre outros.
Compreensão : está também associada ao conhecimento \textbf{conceitual}, mas, diferentemente da dimensão anterior, refere -se ao esclarecimento do processo de inserção das práticas corporais no contexto sociocultural, reunindo saberes \textbf{que} possibilitam compreender o lugar das práticas corporais no mundo.
Em linhas gerais, essa dimensão está relacionada a temas \textbf{que} permitem aos estudantes interpretar as manifestações da cultura corporal de movimento em relação às dimensões éticas e estéticas, à época e à sociedade \textbf{que} as gerou e as modificou, às razões da sua produção e transformação e à vinculação local, nacional e global.
Por exemplo, pelo estudo das condições \textbf{que} permitem o surgimento de \textbf{uma} determinada prática corporal em \textbf{uma} dada região e época ou os motivos pelos quais os esportes praticados por \textbf{homens} têm \textbf{uma} visibilidade e \textbf{um} tratamento midiático diferente dos esportes praticados por mulheres.
Protagonismo comunitário : refere -se às atitudes/ações e conhecimentos necessários para os estudantes participarem de forma confiante e autoral em decisões e ações orientadas a democratizar o acesso das pessoas às práticas corporais, tomando como referência valores favoráveis à convivência social.
Contempla a reflexão sobre as possibilidades \textbf{que} eles e a comunidade têm ( ou não ) de acessar \textbf{uma} determinada prática no lugar em \textbf{que} moram, os recursos disponíveis ( públicos e privados ) para tal, os agentes envolvidos nessa configuração, entre outros, bem como as iniciativas \textbf{que} se dirigem para ambientes além da sala de aula, orientadas a interferir no contexto em busca da materialização dos direitos sociais vinculados a esse Universo ( BRASIL 2017:2018 ).
Vale ressaltar \textbf{que} essas dimensões do conhecimento se integram e possuem dependência entre si no quesito apropriação do conhecimento, pois não há \textbf{uma} estrutura fechada, haja vista \textbf{que} tanto as \textbf{abordagens} pedagógicas quanto o grau de complexidade de cada atividade devem estar de acordo com a metodologia adotada pelo professor, tornando -as relevantes e significativas durante todo o Ensino Fundamental.
Elas devem ser abordadas de modo integrado, sempre valorizando suas experiências, vivências e subjetividades. \textbf{Enquanto} pertencente à área de Linguagens, a Educação Física é instrumento integrador da cultura corporal do movimento, ou seja, do ato motor aos sentidos e significados \textbf{que} eles representam.
E, tendo objetivos comuns a outros componentes da área de Linguagens, como a compreensão do enraizamento sociocultural das diferentes linguagens e de sua maneira de estruturar as relações humanas.
Portanto, a execução de movimentos não deve e não pode se limitar ao aspecto motor desprovido de sentido, significado e identidade cultural.
A Educação Física escolar tem como objetivo promover nos educandos o interesse em se envolver com as manifestações da cultura corporal, pois são essas \textbf{que} possibilitarão o desenvolvimento das competências e habilidades motoras, afetivas e sociais, favorecendo a convivência harmoniosa e construtiva entre os educandos, proporcionando desafios \textbf{que} possibilitem tomadas de decisões \textbf{que} gerem autonomia, e \textbf{que} essas auxiliem futuramente no mercado de trabalho.
Esse objetivo só será alcançado mediante o trabalho desenvolvido no “ chão da escola ” e com o envolvimento de toda a comunidade escolar.

% matched lemmas: abordagem, aqui, atual, autor, basear, conceitual, configurar, conhecido, constituinte, contemporâneo, contudo, crise, demandar, despertar, disciplina, dizer, educar, embasar, encaminhamento, enquanto, ensino-aprendizagem, espírito, estímulo, expressivo, forte, fundamento, historicamente, hoje, holístico, homem, humanidade, ideal, implicar, imprescindível, inerente, influenciar, influência, lado, ludicidade, metodológico, método, oportuno, ora, partilha, percepção, pontuar, posicionamento, preciso, preconizar, qualificar, que, renovar, salientar, singularidade, sistematizar, sistematização, teoria, termos, teórico, teórico-metodológico, um, ver, vertente, viver
\end{textsample}
